\section{Implementation}
\label{sec:impl}

\sysname is written in C++23 and built with \texttt{cmake}. The three
subsystems this paper exercises---the capability manager (\Mone), the
IPC ring (\Mthree), and the M-21 assembly primitives---come to about
7,400 lines of source.

\subsection{Capability manager (\Mone)}
\label{sec:impl:cap}

\Mone is one C++ class, \texttt{CapabilityManager}, holding a
fixed-capacity slot pool sized at construction
(\texttt{DEFAULT\_MAX\_TOKENS}\,=\,4096). The pool is a flat
heap-allocated array of \texttt{TokenSlot}s; neither validate nor
revoke allocates.

\para{Concurrency}
\texttt{validate()} is lock-free: an acquire load of the global
\texttt{m\_uCurrentEpoch}, the $O(1)$ slot lookup, done.
\texttt{create()}, \texttt{derive()}, and \texttt{revoke()} serialise
through a small spinlock (\texttt{m\_bSpinlock}) that guards slot
allocation, the descendant search, and the bookkeeping fields.
Releasing the lock is a release store of the lock byte, which gives
any later \texttt{validate()} on another core the synchronizes-with
edge it needs to observe the new state.

\texttt{revoke()} publishes deactivation in a fixed order: revoke-epoch
stamp (release store), active-flag clear (release store), then secure
wipe of the token bytes. A \texttt{validate()} racing the revoke
either observes one of the two flags, or---having loaded
\texttt{active==true} just before the clear---fails the constant-time
compare against already-wiped bytes. There is no interleaving in
which a revoked token validates.

\para{Token-id minting}
\texttt{generateTokenId(slot, out)} fills the identifier from RDRAND
(via the M-21 \texttt{asm\_rdrand64} wrapper) and packs the slot
index into the high 32 bits of \texttt{out[0]}. Both 64-bit halves
are checked for RDRAND failure. If RDRAND does fail---which should
not happen on any supported part---we fall back to an
epoch-XOR-constant value and log a warning; graceful degradation
beats minting nothing.

\subsection{Ring buffer (\Mthree)}
\label{sec:impl:ring}

\Mthree is a single-producer, single-consumer ring with an
MPSC-capable claim path. Ring data and a three-cache-line control
block live in a \texttt{memfd\_create} region mapped
\texttt{MAP\_SHARED} into both endpoints. The layout
(\Cref{fig:cb}) keeps the claim/commit indices and the read index on
separate cache lines, so producers and the consumer do not
false-share.

\begin{figure}[t]
\centering\small
\begin{verbatim}
+--------------------------------+ line 0
| atomic<u64> write_index        |
| atomic<u64> write_claim_index  |
+--------------------------------+ line 1
| atomic<u64> read_index         |
+--------------------------------+ line 2
| u32 channel_id, ring_bytes ... |
+--------------------------------+ 192B total
\end{verbatim}
\caption{\texttt{ChannelControlBlock} layout. A static\_assert pins
the struct at exactly 3$\times$64\,B.}
\label{fig:cb}
\end{figure}

\para{Two-phase claim/commit}
Payloads up to 256\,B claim their byte range with one
\texttt{fetch\_add} on \texttt{write\_claim\_index}---wait-free.
Larger payloads CAS with bounded retries. After writing its data, the
producer advances \texttt{write\_index} with a CAS that spins on the
previously committed value, so readers always observe one contiguous,
ordered byte stream no matter how many producers raced.

\para{The capability gate}
\texttt{writeGated(\&\_gate, data, len)} calls
\texttt{gate.m\_pManager->validate(gate.m\_token, IPC\_SEND)}; on
failure it returns \texttt{PERMISSION\_DENIED} with
\texttt{write\_claim\_index} untouched, and on success it forwards to
\texttt{write(data, len)}. The gate sits in front of the claim
precisely so a revoked token cannot consume claim space.

\para{The blocking variant}
\texttt{readWaitGated()} parks on \texttt{write\_index.wait()}---a
futex on Linux---and re-validates after each wake. That extra
validate per wake is what keeps revocation airtight across sleeping
readers.

\subsection{M-21 constant-time primitives}
\label{sec:impl:asm}

\Mtwentyone ships seven routines hand-written in NASM against the
System~V AMD64 ABI: \texttt{asm\_ct\_compare} (constant-time 16-byte
compare), \texttt{asm\_secure\_wipe} (a memset the optimiser cannot
delete), \texttt{asm\_rdrand64}, \texttt{asm\_lfence},
\texttt{asm\_cache\_flush}, \texttt{asm\_spectre\_barrier}, and
\texttt{asm\_stack\_guard}.

None of these are contributions---libsodium and BoringSSL ship
equivalents---but the constant-time discipline matters because
\texttt{asm\_ct\_compare} sits on the IPC fastpath. They appear here
rather than in \Cref{sec:design} because the design is agnostic to
which comparator implements it.

\subsection{CBMC monotonicity proof}
\label{sec:impl:cbmc}

We model \texttt{CapabilityManager}'s permission arithmetic in a
small flat-C translation
(\texttt{formal/cbmc/capability\_monotonicity.c}, 280 LOC) and prove
three properties with CBMC~\cite{cbmc}:

\begin{description}
\itemsep0pt
\item[Subset-on-success] If \texttt{derive(parent, child\_perms)}
succeeds, then \texttt{(child\_perms \& \textasciitilde
parent\_perms) == 0}.
\item[Rejection-bijection] \texttt{derive} fails exactly when
\texttt{(child\_perms \& \textasciitilde parent\_perms) != 0}.
\item[No-recovery-through-chain] Along any derivation chain
$t_0 \to t_1 \to \dots \to t_n$, the permissions of $t_n$ are a
subset of those of $t_0$.
\end{description}

CBMC reports 0 of 37 checks failing under
\texttt{formal/cbmc/Makefile}'s \texttt{verify} target. We are
careful about what this is: \textit{spec correspondence on the one
property the runtime must never get wrong}. It is not functional
verification in the seL4 sense~\cite{klein2009sel4}, and we do not
present it as such.

\subsection{Reproducibility}
\label{sec:impl:repro}

The artefact---source, baselines, sweep driver, plotting---is
available at \ifanon\textcolor{red}{[anonymous URL during review]}\else
\url{https://github.com/KernelArch-Lab/Astra}\fi.
One command, \texttt{scripts/run\_paper1\_sweep.sh}, builds every
harness, runs each in the same TSC domain on the same machine,
merges the repetitions into \texttt{artefact/paper1\_figure\_1.csv},
and drives the figure and table generators.
\Cref{tab:artefact} maps every figure and claim to its source.

\begin{table}[t]
\centering\scriptsize
\setlength{\tabcolsep}{2pt}
\begin{tabular}{lll}
\toprule
Figure / claim         & Source script         & Output \\
\midrule
Fig 1 latency          & \texttt{sweep.sh} \tiny\texttt{+plot 1} & \texttt{figure\_1.pdf} \\
Fig 2 O(1) validate    & \texttt{bench\_pool\_scaling}    \tiny\texttt{+plot 2} & \texttt{figure\_2.pdf} \\
Tbl 1 256\,B comparison & \texttt{plot\_paper1\_figure.py}                       & \texttt{table\_1.tex}  \\
Throughput             & \texttt{bench\_throughput}                              & CSV                          \\
MPSC scaling           & \texttt{bench\_throughput\_mpsc}                        & CSV                          \\
Revoke window          & \texttt{bench\_revocation\_latency}                     & CSV                          \\
Cycles / cache misses  & \texttt{bench\_perfcounters}                            & CSV                          \\
Monotonicity proof     & \texttt{make -C formal/cbmc verify}                     & CBMC report                  \\
Concurrency stress     & \texttt{test\_capability\_concurrency}                  & ctest pass                   \\
Survivability          & \texttt{test\_capability\_survivability}                & ctest pass                   \\
Gate concurrency       & \texttt{test\_sprint5\_gate\_concurrency}               & ctest pass                   \\
Property test          & \texttt{test\_capability\_property}                     & ctest pass                   \\
\bottomrule
\end{tabular}
\caption{Artefact map (\texttt{sweep.sh} abbreviates
\texttt{scripts/run\_paper1\_sweep.sh}). Every numerical and
functional claim in the paper has a runnable script behind it.}
\label{tab:artefact}
\end{table}
